%% ============================================================
%% Chapter 4 — System Architecture
%% ============================================================
\chapter{System Architecture}
\label{chap:architecture}

\section{Overview}
\label{sec:arch_overview}

The end-to-end pipeline consists of five conceptual stages that transform a raw
question into either a final answer or an abstention signal. Figure~\ref{fig:architecture}
shows the full pipeline as a TikZ diagram.

\begin{enumerate}
  \item \textbf{Data Ingestion \& Preprocessing.}  The SQuAD~2.0 JSON file is
        loaded and parsed. Each article's passages and questions are extracted.
        An article-level stratified split divides the data into train, validation,
        and test partitions, ensuring no article straddles two splits.

  \item \textbf{BM25 Indexing.}  For each split, an Okapi BM25 index is built
        over the split's passages. The index stores inverted lists over
        lower-cased, whitespace-tokenised terms, weighted by IDF.

  \item \textbf{Retrieval \& Feature Extraction.}  At query time, the question
        is used to query the BM25 index. The top-5 scored passages are
        retrieved; five retrieval features are computed from their scores
        (Table~\ref{tab:features}).

  \item \textbf{QA Reader.}  The top-1 retrieved passage is concatenated with
        the question and fed to the \texttt{deepset/roberta-base-squad2}
        reader. The reader produces a predicted answer span together with a
        \texttt{qa\_score} and a \texttt{qa\_no\_answer\_probability}. Six
        additional reader and query features are derived.

  \item \textbf{Abstention Decision.}  One of the three strategies
        (Always Answer, BM25 Threshold, Answerability Classifier) processes the
        combined feature vector and returns either the predicted answer or an
        abstention token~$\varnothing$.
\end{enumerate}

\section{TikZ Pipeline Diagram}
\label{sec:arch_diagram}

\begin{figure}[htbp]
  \centering
  \begin{tikzpicture}[
      node distance = 1.1cm and 1.6cm,
      >=Stealth,
      block/.style = {rectangle, draw, rounded corners=4pt,
                      fill=blue!8, text width=3.2cm,
                      minimum height=1.1cm, align=center,
                      font=\small},
      decision/.style = {diamond, draw, fill=orange!15,
                         text width=2.2cm, align=center,
                         inner sep=1pt, font=\small},
      io/.style     = {rectangle, draw, fill=green!10,
                       text width=2.5cm, minimum height=1cm,
                       align=center, font=\small},
      result/.style = {rectangle, draw, fill=red!10,
                       text width=2.5cm, minimum height=1cm,
                       align=center, font=\small},
      line/.style   = {draw, -Stealth, thick},
    ]

    %% Nodes
    \node[io]       (input)    {Question $q$};
    \node[block,    right=of input]  (bm25idx)  {BM25\\Index};
    \node[block,    right=of bm25idx](retrieval){Top-5 Passage\\Retrieval};
    \node[block,    right=of retrieval](reader){RoBERTa\\Reader};

    %% Feature extraction below retrieval
    \node[block,    below=1.5cm of retrieval](features){Feature\\Extraction\\(11 features)};

    %% Abstention decision
    \node[decision, below=1.5cm of reader](abstain){Abstention\\Strategy};

    %% Outputs
    \node[result,   right=1.4cm of abstain](answer){Answer $\hat{a}$};
    \node[result,   below=0.9cm of abstain](abstain_out){Abstain $\varnothing$};

    %% Corpus
    \node[io, left=of bm25idx](corpus){Passage\\Corpus};

    %% Arrows
    \draw[line] (corpus)    -- (bm25idx);
    \draw[line] (bm25idx)   -- (retrieval);
    \draw[line] (input)     |- (retrieval);
    \draw[line] (retrieval) -- (reader);
    \draw[line] (retrieval) -- (features);
    \draw[line] (reader)    |- (features);
    \draw[line] (reader)    -- (abstain);
    \draw[line] (features)  -| (abstain);
    \draw[line] (abstain)   -- node[above]{\small Answer} (answer);
    \draw[line] (abstain)   -- node[right]{\small Abstain} (abstain_out);

  \end{tikzpicture}
  \caption{End-to-end RAG-with-abstention pipeline. The question is used to
           retrieve passages via BM25, the top-1 passage is passed to the
           RoBERTa reader, and the combined 11-dimensional feature vector is
           evaluated by the abstention strategy to decide whether to answer or
           abstain.}
  \label{fig:architecture}
\end{figure}

\section{Component Interactions and Data Flow}
\label{sec:arch_dataflow}

The pipeline is implemented in Python using the \texttt{rank\_bm25} library for
retrieval and the Hugging Face \texttt{transformers} library for the QA reader.
Feature vectors are stored as \texttt{pandas} DataFrames and serialised to CSV
files after each pipeline stage, enabling reproducibility and offline analysis.

\paragraph{BM25 $\rightarrow$ Features.}
After retrieval, the \texttt{top\_bm25\_score}, \texttt{second\_bm25\_score},
\texttt{bm25\_score\_gap}, \texttt{mean\_topk\_bm25}, and
\texttt{source\_passage\_retrieved} features are computed directly from the
ranked list of (score, passage) pairs. These five values, together with the
question metadata features (\texttt{question\_length}, \texttt{question\_token\_count}),
are available before the reader is invoked.

\paragraph{Reader $\rightarrow$ Features.}
The reader's outputs (\texttt{qa\_score}, \texttt{qa\_no\_answer\_probability},
\texttt{predicted\_answer\_length}, \texttt{answer\_nonempty}) are appended to
the feature vector after the reader processes the top-1 passage.

\paragraph{Classifier Training.}
During training, the full 11-feature vectors are labelled with the ground-truth
answerability flag and used to fit the logistic regression classifier. The
trained model and its fitted \texttt{StandardScaler}+\texttt{SimpleImputer}
pre-processors are saved to disk and loaded at test time.

\paragraph{Threshold Selection.}
Both the BM25 threshold $\tau$ and the classifier threshold $\delta$ are
selected by sweeping candidate values on the validation set, measuring F1, and
picking the maximiser. The best classifier threshold found is $\delta^* = 0.28$.
Only then are the thresholds applied to the held-out test set.

\section{Design Principles}
\label{sec:arch_principles}

Several design principles guided the architecture:

\begin{description}
  \item[Modularity.] Each stage produces a well-defined output artifact (CSV
        file or pickled model), allowing stages to be run independently and
        results to be reproduced without re-running the entire pipeline.

  \item[Leakage prevention.] BM25 indices are built per-split; the classifier
        is trained only on training features; thresholds are tuned only on
        validation features. The test set remains entirely held out.

  \item[Interpretability.] The choice of logistic regression over a black-box
        model is deliberate: the learned coefficients can be inspected to
        understand which features most influence the abstention decision.

  \item[Scalability considerations.] While BM25 is efficient at corpus scale,
        the architecture is designed so that the retrieval backbone can be
        swapped for a dense retriever (e.g., DPR) without changing the
        downstream feature extraction or classifier code.
\end{description}
